\section{Functional Composition of the Command Pipeline}

\subsection{Formal Definitions}

Let:
\begin{itemize}
    \item $C$ — the set of all \textbf{commands}.
    \item $R$ — the set of all \textbf{results}.
    \item $E$ — the set of all \textbf{exceptions}.
    \item $\mathcal{S}$ — the set of all \textbf{suppliers}, where $s \in \mathcal{S}$ is a function $s: \varnothing \to R \cup E$.
\end{itemize}

\subsection{Core Mappings}
\[
\begin{aligned}
\text{route} &: C \to H && \text{(selects the appropriate handler)} \\
\text{handle}_h &: C \to R \cup E && \text{(handler processes the command)} \\
\text{exec} &: C \to \mathcal{S} && \text{(executor creates deferred computation)} \\
\text{pre}_m &: C \to C && \text{(middleware transformation)} \\
\text{post}_p &: C \times (R \cup \{\bot\}) \times (E \cup \{\bot\}) \to (R \cup \{\bot\}) \times (E \cup \{\bot\}) && \text{(post-processor)}
\end{aligned}
\]

\subsection{Composed Pipeline Function}
\[
\mathrm{pipeline} = \mathrm{wrapPost} \circ \mathrm{exec}
\]

For a command $c \in C$:
\[
s = \mathrm{exec}(c)
\]
and
\[
\mathrm{wrapPost}(s)() =
\mathrm{finalize}\big(
    \mathrm{fold}_{p \in P}(\mathrm{post}_p,\, \mathrm{tryExec}(s))
\big)
\]

\subsection{Execution Semantics}
\[
\mathrm{tryExec}(s) =
\begin{cases}
    (s(), \bot) & \text{if handler succeeds}\\[4pt]
    (\bot, e) & \text{if handler throws exception } e
\end{cases}
\]

\subsection{Temporal Property}
A more precise formulation of the temporal property, separating orchestration from execution, is:
\begin{align}
\forall c \in C,\, \exists! r \in R, e \in E : \quad
& \underbrace{\mathrm{pre}_{m_1} \circ \dots \circ \mathrm{pre}_{m_n}}_{\text{middleware orchestration}}(c) \notag\\
& \;\to\; 
\underbrace{\mathrm{exec}_h(c)}_{\text{command execution}} \notag\\
& \;\to\; 
\underbrace{\mathrm{post}_{p_1} \circ \dots \circ \mathrm{post}_{p_m}}_{\text{post-processing orchestration}}(c,r,e)
\end{align}

Where:
\begin{itemize}
    \item $\mathrm{pre}_{m_i}$ – middleware transformation and side-effects (validation, logging, etc.) before command execution.
    \item $\mathrm{exec}_h$ – the actual handler invocation performing the command's business logic, producing $r \in R$ or $e \in E$.
    \item $\mathrm{post}_{p_j}$ – post-processors applied sequentially after successful or failed execution for auditing, retries, notifications, etc.
\end{itemize}

This separation makes it explicit that \textbf{pipeline orchestration (middleware and post-processing sequencing) is distinct from the core command execution}, even though the temporal property ensures deterministic and traceable sequencing across the full lifecycle.

\paragraph{Resilience Extension.}
In the presence of systemic failures (e.g., transient network or database unavailability),
the temporal property is extended with fallback semantics.
If $\mathrm{exec}(c)$ cannot complete successfully, a resilience mechanism 
(such as Resilience4J retries) attempts recovery:
\[
\forall c \in C,\ \exists! p \in P :
\big(
    \mathrm{exec}(c) \prec \mathrm{post}_p(c,r,e)
\big)
\ \lor\
\big(
    \mathrm{fail}(c) \Rightarrow \Diamond\, \mathrm{retry}(c)
\big)
\]

If recovery eventually fails, the system propagates an explicit exception,
preserving state consistency and traceability throughout the command lifecycle.


